% Section 7.7, from lectures/L07/appendix/b_exercises.md.

\section{Exercises}
\label{c7:sec:exercises}

\begin{aside}
\textbf{How to work these.} Each exercise is tagged with the kinds of work it asks for:
\kinds{Paper} by hand, \kinds{CircuitVerse} in the browser, \kinds{VHDL} in a module of your own.

\textbf{How to check your work.} Every exercise below that asks you to write a VHDL module ships a
self-checking testbench under \file{lectures/L07/exercises} in the companion repository. Write your
module in its \file{exercises/<module>} directory, using the entity name and the \textbf{port
order} the exercise specifies, then run it with GHDL; \secref{c2:app:c} has the three commands.

From this chapter on, designs reuse modules you wrote earlier, and those are not shipped a second
time: each exercise says which of your own files to copy in and prints the \code{ghdl -a} line that
needs them.

No FPGA board is needed for any exercise. The Quartus synthesis and board-programming steps are
demonstrated in the course, and Appendix~\ref{app:quartus} shows them for a reader with a board;
your job is to get the VHDL right, and the testbench is how you confirm it. Short answers to the
parts you work on paper are in Appendix~\ref{app:answers}; reference solutions to the VHDL are
published in the companion repository, under \file{lectures/L07/solutions}. Try each exercise
yourself before looking at either.
\end{aside}

% ------------------------------------------------------------------------------------------
\exerciseset{Timers by Hand}
\label{c7:set:1}

\exercise{Retargeting the timer by hand}{Paper, CircuitVerse}
\label{c7:ex:1}
Build a second timer by hand in CircuitVerse, following \secref{c7:a2}, but with a target count of
\code{TICK_COUNT = 12} instead of the chapter's \code{10}.

\begin{parts}
  \item Retarget the comparator:
    \begin{itemize}
      \item Write \code{12} in 4-bit binary.
      \item The four XNOR gates stay exactly where they are. Say which of their constant inputs
        change, and to what.
      \item State, in one sentence, why an equality comparator built from XNORs needs no gate
        added, removed or rewired to compare against a different value.
    \end{itemize}
  \item Build it:
    \begin{itemize}
      \item Reuse your 4-bit counter and comparator from \secref{c7:a2}, changing only the
        constants.
      \item Set the clock period to \code{1000 ms} and confirm \code{timeout} fires once every
        thirteen clock edges, for exactly one edge each time.
    \end{itemize}
  \item Now remove the clear, so that \code{timeout} no longer forces the counter back to
    \code{0000}:
    \begin{itemize}
      \item Predict, before running it, how often \code{timeout} fires now.
      \item Run it and confirm.
      \item Explain the result in terms of the counter's own wraparound (\secref{c6:a1}).
    \end{itemize}
  \item Deliberately break the comparator: disconnect one XNOR's output from the AND gate and tie
    that AND input high instead.
    \begin{itemize}
      \item Predict which counts now assert \code{timeout}, before simulating.
      \item Confirm by stepping the counter through a full cycle.
      \item This is what ``three of the four bits match'' looks like, and it is the most common way
        a hand-drawn comparator goes wrong.
    \end{itemize}
\end{parts}
\note[Tip] counting \code{0} through \code{12} inclusive is thirteen edges, not twelve.
\Secref{c7:a1}'s note on \code{TICK_COUNT + 1} is the same off-by-one, and it is worth getting
straight here, where you can count the edges by eye, rather than later at 50\,MHz.

\exercise{Reading the timer module}{Paper}
\label{c7:ex:2}
Answer each of the following in one or two sentences, referring to the \code{timer} module in
\secref{c7:a3}:

\begin{parts}
  \item Why is \code{timeout} a single-cycle pulse, rather than a signal that stays high once the
    target count is reached?
  \item While a running timer's \code{enable} input is held low, what happens to its internal
    counter from one clock edge to the next: does it reset, pause, or keep counting?
  \item Why does the \code{timer} take \code{reset_s2_n} (the synchronized reset) rather than the
    raw \code{reset_n} as its reset input?
  \item The CircuitVerse timer in Exercise~\ref{c7:ex:1} needed rewiring to change its target
    count, while the VHDL module needs no edit at all. What replaces the rewiring, and at what point
    is its value fixed?
\end{parts}

% ------------------------------------------------------------------------------------------
\exerciseset{Timers in VHDL}
\label{c7:set:2}

\exercise{Tick counts at 50\,MHz}{Paper}
\label{c7:ex:3}
The DE0-CV board's system clock is \code{50 MHz}. For each target below, calculate the
\code{TICK_COUNT} generic you would pass to the \code{timer} module from \secref{c7:a3}:

\begin{parts}
  \item A \code{timeout} pulse once per second.
  \item A \code{timeout} pulse every \code{250} ms.
  \item A \code{timeout} pulse twenty times per second (\code{20 Hz}).
\end{parts}
\note[Tip] The relationship between a period and its tick count is:
\[
  \bool{TICK\_COUNT} = \bool{seconds} \times 50{,}000{,}000 - 1
\]
The \code{- 1} is the same off-by-one Exercise~\ref{c7:ex:1} turned on: the counter runs \code{0}
through \code{TICK_COUNT} inclusive, so a period is \code{TICK_COUNT + 1} clock cycles, not
\code{TICK_COUNT}.

\exercise{Writing the timer}{VHDL, Paper}
\label{c7:ex:4}
Write the \code{timer} module itself, as an entity named \code{timer}.

The course builds it live and \secref{c7:a3} prints it in full, so this is not a puzzle. Write it
anyway, without copying: every design from here to the end of the book instantiates this module,
including both capstones, and it is worth having built the thing you are about to reuse six times.

The entity has one generic:

\begin{booktable}{@{}p{5.2em}p{3.8em}p{5.4em}L@{}}
  \thead{Generic} & \thead{Type} & \thead{Default} & \thead{Description} \\
  \midrule
  \code{TICK_COUNT} & \code{natural} & \code{50_000_000} & Ticks to count before a timeout. One
    second at the DE0-CV's \code{50 MHz} clock, rounded: by Exercise~\ref{c7:ex:3}'s rule the
    exact value is \code{49_999_999}, and 20\,ns of error in a second is not worth the less
    readable constant. \\
\end{booktable}

\noindent and these ports:

\begin{booktable}{@{}p{5.6em}p{4.8em}p{4.6em}L@{}}
  \thead{Port} & \thead{Direction} & \thead{Type} & \thead{Description} \\
  \midrule
  \code{clock} & in & \code{std_logic} & System clock. \\
  \code{reset_s2_n} & in & \code{std_logic} & Active-low, already-synchronized reset.
    Asynchronous: it clears the counter and \code{timeout} with no clock edge involved. \\
  \code{enable} & in & \code{std_logic} & Timer enable. While low the counter \textbf{holds} its
    value, and resumes from there rather than restarting. \\
  \code{timeout} & out & \code{std_logic} & One-cycle pulse, asserted every
    \code{TICK_COUNT + 1} cycles while enabled. \\
\end{booktable}

\noindent The architecture must:
\begin{itemize}
  \item Declare the internal counter as \code{natural range 0 to TICK_COUNT}, so the register is
    sized by the generic rather than left unconstrained.
  \item Use one process, sensitive to \code{clock} and \code{reset_s2_n}, following the
    clocked-process template from \secref{c3:a6}.
  \item Clear both the counter and \code{timeout} on reset, outside the clocked branch.
  \item Drive \code{timeout} low unconditionally at the top of the clocked branch, and high only on
    the edge that clears the counter. That default-then-override is what makes the pulse exactly
    one cycle wide rather than two.
\end{itemize}
Then answer, in a sentence each:
\begin{itemize}
  \item \code{enable = '0'} must \textbf{pause} the counter, not clear it. Suppose you cleared it
    instead. Predict what the reference testbench reports, and which of the two properties above it
    is measuring when it does. Then run it and see. (Chapter~\ref{c8:ch}'s \code{fsm_led} depends
    on the pause, which is why the testbench bothers.)
  \item The counter counts \code{0} through \code{TICK_COUNT} inclusive, so a period is
    \code{TICK_COUNT + 1} cycles. Where in your code is that ``+ 1'' actually written? It is not a
    literal anywhere.
  \item \Secref{c7:a2}'s hand-drawn \code{timeout} is a combinational decode of the count; yours is
    registered. Both have the same period. Which cycle does each one fire on, and why does the
    difference matter for anything downstream that samples \code{timeout}?
\end{itemize}

\modulefigure{timer}

\selfcheck name your entity \code{timer}, with generic \code{TICK_COUNT} (\code{natural}, default
\code{50_000_000}), inputs \code{clock}, \code{reset_s2_n}, \code{enable}, and output
\code{timeout}, declared in that order; its testbench is in \file{exercises/timer}. It runs three
full periods, so a one-shot passes nothing, and pauses mid-count to confirm the count resumed rather
than restarted.

\textbf{Keep this file.} Exercise~\ref{c7:ex:5} below reuses it, so does \code{walking_led} in
Exercise~\ref{c7:ex:6}, and so do both Chapter~\ref{c8:ch} capstones. Every later exercise that
needs a timer asks you to copy this one in.

\exercise{A one-second blinker}{VHDL, Paper}
\label{c7:ex:5}
Write an entity named \code{blinker} that blinks an LED on and off once per second, by reusing the
\code{timer} module you just wrote rather than counting clock cycles itself.

The entity has one generic:

\begin{booktable}{@{}p{5.2em}p{3.8em}p{5.4em}L@{}}
  \thead{Generic} & \thead{Type} & \thead{Default} & \thead{Description} \\
  \midrule
  \code{TICK_COUNT} & \code{natural} & \code{25_000_000} & A half-second at the DE0-CV's
    \code{50 MHz} clock. Expose it as a generic rather than hard-coding the count, so a testbench
    can shorten the blink; the FPGA build just uses the default. \\
\end{booktable}

\noindent and these ports:

\begin{booktable}{@{}p{5.6em}p{4.8em}p{4.6em}L@{}}
  \thead{Port} & \thead{Direction} & \thead{Type} & \thead{Description} \\
  \midrule
  \code{clock} & in & \code{std_logic} & System clock. \\
  \code{reset_n} & in & \code{std_logic} & Active-low reset, raw from the outside world.
    Asynchronous: assert it and \code{led} goes out immediately, without waiting for a clock
    edge. \\
  \code{led} & out & \code{std_logic} & Cleared to \code{'0'} while the reset is asserted, and left
    there until the first timeout after the reset is released. \\
\end{booktable}

\noindent Note that this entity takes the \textbf{raw} \code{reset_n} rather than an
already-synchronized \code{reset_s2_n}. Synchronizing it is part of the exercise, and it is what
\secref{c4:a5} means by ``exactly one module in a design is the exception'': here that module is
\code{blinker} itself.

The architecture must:
\begin{itemize}
  \item Instantiate your \code{reset_sync} from Chapter~4's Exercise~\ref{c4:ex:8}, and use its
    \code{reset_s2_n} output as the reset for everything inside.
  \item Instantiate the \code{timer} module from \secref{c7:a3}:
    \begin{itemize}
      \item Pass your own \code{TICK_COUNT} generic straight through to the timer's
        \code{TICK_COUNT}.
      \item Hold its \code{enable} input high.
    \end{itemize}
  \item Toggle \code{led} in a synchronous process whenever the timer's \code{timeout} pulse fires.
\end{itemize}
Do not reimplement the counting logic, or the synchronizer, yourself.

Instead:
\begin{itemize}
  \item Reuse \code{timer} and \code{reset_sync} unchanged, via a \code{generic map} and
    \code{port map}, following the composition pattern in \secref{c7:a4}.
  \item Explain why a \textbf{half}-second timer period produces a \textbf{one}-second blink cycle
    (one full on-then-off period).
\end{itemize}
Then answer, in a sentence each:
\begin{itemize}
  \item The LED must be off during reset, and the testbench checks that it goes off with no clock
    edge in between. Which half of the assert-async-release-sync pattern does that check, and which
    half does it not?
  \item Your \code{timer} is reset by \code{reset_s2_n}, two edges after \code{reset_n} is
    released. What would go wrong if you fed it the raw \code{reset_n} instead? Name the failure,
    not just the rule.
\end{itemize}

\modulefigure{blinker}

\selfcheck name your entity \code{blinker}, with generic \code{TICK_COUNT} (\code{natural}, default
\code{25_000_000}), inputs \code{clock}, \code{reset_n}, and output \code{led}, declared in that
order; its testbench is in \file{exercises/blinker}. \textbf{Copy your own \file{reset_sync.vhd}
and \file{timer.vhd} into that directory} before you build; neither is shipped there, because you
wrote both. It overrides \code{TICK_COUNT} with a small value and checks that the LED toggles
repeatedly, holds for exactly one timer period between toggles, and that the first toggle arrives
late by your synchronizer's two edges (see \secref{c2:app:c}):

\begin{shell}
cd lectures/L07/exercises/blinker
cp ../../../L04/exercises/reset_sync/reset_sync.vhd .   # the two you wrote yourself
cp ../timer/timer.vhd .
ghdl -a --std=93 reset_sync.vhd timer.vhd blinker.vhd blinker_tb.vhd
ghdl -e --std=93 blinker_tb
ghdl -r --std=93 blinker_tb --assert-level=error --stop-time=10ms
\end{shell}

% ------------------------------------------------------------------------------------------
\exerciseset{The Walking LED}
\label{c7:set:3}

A single lit LED walking along a row, one position per timer tick, started and stopped with a
button. The course builds it live, and \secref{c7:a5} writes it up.

This is the payoff of the last four chapters, and it writes almost no logic of its own: three
modules, two of them carried unchanged from Chapter~\ref{c4:ch} and one written in this chapter,
composed. Exercise~\ref{c7:ex:6} builds it, and Exercises~\ref{c7:ex:7} and~\ref{c7:ex:8} modify
what you built.

\exercise{The walking LED}{VHDL, Paper}
\label{c7:ex:6}
Write an entity named \code{walking_led}.

Like the \code{timer} of Exercise~\ref{c7:ex:4}, the course builds it live and \secref{c7:a5}
prints it in full, so this is not a puzzle either. Write it anyway, without copying: composing
three modules you wrote yourself is the whole point of the exercise, and the capstones in
Chapter~\ref{c8:ch} ask for it again at twice the size.

The entity has two generics:

\begin{booktable}{@{}p{5.2em}p{3.8em}p{5.4em}L@{}}
  \thead{Generic} & \thead{Type} & \thead{Default} & \thead{Description} \\
  \midrule
  \code{LED_COUNT} & \code{natural} & \code{8} & How many LEDs the bit walks along. The row is this
    wide, and the walk wraps after this many steps. \\
  \code{TICK_COUNT} & \code{natural} & \code{25_000_000} & Passed straight through to your
    \code{timer}, so the bit moves one position per timer period: a half-second step at the
    DE0-CV's \code{50 MHz} clock. \\
\end{booktable}

\noindent and these ports:

\begin{booktable}{@{}p{4.6em}p{4.8em}p{10.8em}L@{}}
  \thead{Port} & \thead{Direction} & \thead{Type} & \thead{Description} \\
  \midrule
  \code{clock} & in & \code{std_logic} & System clock. \\
  \code{reset_n} & in & \code{std_logic} & Active-low reset, raw from the outside world.
    Synchronizing it is your job, as it was in \code{blinker}. \\
  \code{button_n} & in & \code{std_logic} & Active-low push button, raw. Each press starts or stops
    the walk. \\
  \code{led} & out & \code{std_logic_vector(}\allowbreak\code{LED_COUNT-1 downto 0)} & The LED
    row. Exactly one bit is lit at any time; on reset it is bit \code{0}. \\
\end{booktable}

\noindent The architecture composes three modules you have already written, and adds two
processes:
\begin{itemize}
  \item Instantiate your \textbf{\code{reset_sync}}, and use its \code{reset_s2_n} as the reset for
    everything else in the design, including the two instances below.
  \item Instantiate your \textbf{\code{button_sync}} with \code{generic map(1)}, for a one-cycle
    pulse per press. Its button ports are \code{COUNT}-bit vectors while \code{button_n} here is a
    scalar, so bridge the width with a one-element \code{std_logic_vector(0 downto 0)} in each
    direction. \Secref{c7:a5} shows this.
  \item Instantiate your \textbf{\code{timer}} with your own \code{TICK_COUNT}, driving its
    \code{enable} from the shift-enable flag below and taking its \code{timeout} as the shift tick.
  \item A process toggling \textbf{\code{shift_enable}} on each button pulse, cleared to \code{'0'}
    on reset. The walk therefore starts stopped.
  \item A process holding the \textbf{shift register} that drives \code{led}: on reset, load a
    single lit bit at position \code{0} and clear the rest; on each timer \code{timeout}, rotate one
    step toward the MSB, wrapping the top bit back around into bit \code{0}. That is a PISO
    register (\secref{c6:a4}) with its own serial output wired back to its serial input, which is
    what turns a one-shot shift-out into an endless walk.
\end{itemize}
Then answer, in a sentence each:
\begin{itemize}
  \item \code{shift_enable} drives the timer's \code{enable}, so stopping the walk stops the timer
    too. Exercise~\ref{c7:ex:4} made you show that \code{enable} \textbf{pauses} the count rather
    than clearing it. What does a reader see on the board, on the press that restarts the walk, if
    your timer clears instead?
  \item \code{button_sync} is instantiated with \code{COUNT = 1} here and with \code{COUNT = 2} in
    Chapter~\ref{c8:ch}'s \code{fsm_led}, from the same file with no edit. What would you have had
    to write instead if the width were fixed in the module rather than passed as a generic?
\end{itemize}

\modulefigure{walking_led}

\selfcheck name your entity \code{walking_led}, with generics \code{LED_COUNT} (\code{natural},
default \code{8}) and \code{TICK_COUNT} (\code{natural}, default \code{25_000_000}), inputs
\code{clock}, \code{reset_n}, \code{button_n}, and output \code{led}, declared in that order; its
testbench is in \file{exercises/walking_led}. It overrides the generics (\code{LED_COUNT = 4},
\code{TICK_COUNT = 3}) so a full lap takes a handful of clock cycles rather than two seconds, and it
checks the reset pattern, that nothing moves before the first press, and that the shift count over a
fixed window matches the pace the timer sets. Copy in the three modules it composes first:

\begin{shell}
cd lectures/L07/exercises/walking_led
cp ../../../L04/exercises/reset_sync/reset_sync.vhd .   # the three you wrote yourself
cp ../../../L04/exercises/button_sync/button_sync.vhd .
cp ../timer/timer.vhd .
ghdl -a --std=93 reset_sync.vhd button_sync.vhd timer.vhd \
                 walking_led.vhd walking_led_tb.vhd
ghdl -e --std=93 walking_led_tb
ghdl -r --std=93 walking_led_tb --assert-level=error --stop-time=10ms
\end{shell}

\exercise{Changing the walking LED}{VHDL, Paper}
\label{c7:ex:7}
Make the following changes one at a time, re-running the testbench after each. Parts \textbf{a)}
and \textbf{b)} change only the testbench's constants --- that they need no edit to your
\file{walking_led.vhd} at all is the point of them. Part \textbf{c)} is the one that touches the
module.

\begin{parts}
  \item Change the walking speed:
    \begin{itemize}
      \item Edit the \code{TICK_COUNT} constant in the \textbf{testbench} from \code{3} to
        \code{6}.
      \item It should still pass, and this time because the testbench adapts with you: its
        observation window is a fixed \code{PULSES * 6} cycles, and it derives the shift count it
        expects from \code{TICK_COUNT}, one shift per \code{TICK_COUNT + 1} cycles.
      \item Now find where that adaptation stops meaning anything. Keep raising \code{TICK_COUNT}
        past the window length and work out what the expected count becomes, and what the check is
        then actually proving. A test that passes for the wrong reason is worth being able to
        recognize.
      \item Then explain why changing the tick count cannot break the \emph{design}, only slow it
        down.
    \end{itemize}
  \item Change the width:
    \begin{itemize}
      \item Edit the testbench's \code{LED_COUNT} constant from \code{4} to \code{8}.
      \item It should still pass, with no change to your \file{walking_led.vhd} at all. Explain
        what the module would have needed to look like for this to require an edit.
    \end{itemize}
  \item Load two lit bits on opposite sides of the register on reset (e.g.\ \code{"10001000"} for
    \code{LED_COUNT = 8}), and let them walk in lockstep:
    \begin{itemize}
      \item Identify the single line in your \file{walking_led.vhd} you need to change.
      \item Before running anything, \textbf{predict} what the reference testbench will report.
      \item Then run it and see. Explain the result: is the design wrong, or is the testbench
        asserting something that is no longer the specification?
    \end{itemize}
\end{parts}
\note[Tip] Both \code{LED_COUNT} and \code{TICK_COUNT} are generics, so parts \textbf{a)} and
\textbf{b)} need no change to the module's body at all. That is the entire point of a generic: one
module, many sizes, no edits.

\exercise{Reversing the walk}{VHDL, Paper}
\label{c7:ex:8}
The current design only ever walks in one direction (toward the MSB, wrapping the top bit back
around to the bottom). Modify your \file{walking_led.vhd} so that a \textbf{second} button reverses
the walking direction.

\begin{parts}
  \item Add a second, synchronized button input:
    \begin{itemize}
      \item Route it through the same \code{button_sync} instance, passing \code{2} for its
        \code{COUNT} generic (\code{generic map(2)}, positionally, as everywhere else in this
        book).
    \end{itemize}
  \item Add a \code{direction} signal:
    \begin{itemize}
      \item Toggle it on the new button's edge, exactly the way \code{shift_enable} is toggled by
        the first button.
    \end{itemize}
  \item Make the rotation in \code{SHIFT_PROCESS} depend on \code{direction}:
    \begin{itemize}
      \item When \code{direction = '0'}, rotate toward the MSB (the book's convention from
        \secref{c6:a2}, and what you wrote in Exercise~\ref{c7:ex:6}):
\begin{vhdlcode}
shift_reg <= shift_reg(LED_COUNT-2 downto 0) & shift_reg(LED_COUNT-1);
\end{vhdlcode}
      \item When \code{direction = '1'}, rotate the other way, toward bit 0:
\begin{vhdlcode}
shift_reg <= shift_reg(0) & shift_reg(LED_COUNT-1 downto 1);
\end{vhdlcode}
    \end{itemize}
    This is the one place in the book where a shift register deliberately runs against the
    convention: reversing direction is the entire point of the exercise. Note that the
    second expression is exactly the mirror image \secref{c6:a2} warns you to watch for in other
    people's code.
  \item Your entity's ports have now changed: \code{button_n} is a two-bit vector rather than a
    single bit.
    \begin{itemize}
      \item Explain why \file{walking_led_tb.vhd} can no longer be used against your version, and
        what precisely GHDL will complain about when you try.
      \item This is the ordinary cost of changing an interface, and it is worth feeling once: a
        testbench is written against a specific entity, and widening a port breaks every
        instantiation of it, including the ones you did not write.
      \item Describe, in prose, the sequence of presses you would use to convince yourself the
        design works, and what you would expect the LEDs to do at each step. That sequence is
        exactly what a testbench for the new entity would have to drive.
    \end{itemize}
\end{parts}
